\documentclass[11pt,a4paper]{article}
\usepackage[margin=1in]{geometry}
\usepackage[T1]{fontenc}
\usepackage[utf8]{inputenc}
\usepackage{amsmath,amssymb,booktabs,longtable,array}
\usepackage{xcolor}
\usepackage[hidelinks]{hyperref}
\newcolumntype{L}[1]{>{\raggedright\arraybackslash}p{#1}}
\newcommand{\code}[1]{\begingroup\urlstyle{tt}\footnotesize\nolinkurl{#1}\endgroup}
\setlength{\LTleft}{0pt}
\setlength{\LTright}{0pt}
\setlength{\tabcolsep}{4pt}
\setlength{\emergencystretch}{3em}
\allowdisplaybreaks
\renewcommand{\arraystretch}{1.2}

\title{Collision-Based Export Formula Report}
\author{Generated from stress\_collision\_based.py}
\date{}
\begin{document}
\maketitle
\noindent Formulas implemented in \code{stress_collision_based.py}\par
\medskip
\noindent This document summarizes the exact formulas currently implemented in \code{stress_collision_based.py} to produce exported CSV data.

\section*{Output files}
\begin{longtable}{@{}L{0.18\textwidth}L{0.32\textwidth}L{0.42\textwidth}@{}}
\toprule
\textbf{Type} & \textbf{File name pattern} & \textbf{Content}\\
\midrule
\endhead
Summary CSV & \code{<dem>_collision_over_time.csv} & One row per processed timestep.\\
Extracted CSV & \code{<dem>_collision_extracted_at_<time>.csv} & Per-collision values at the selected extract time.\\
\bottomrule
\end{longtable}

\section*{Notation}
\begin{longtable}{@{}L{0.22\textwidth}L{0.70\textwidth}@{}}
\toprule
\textbf{Symbol} & \textbf{Meaning}\\
\midrule
\endhead
$k$ & Current timestep index.\\
$t_k$ & Current sample time returned by EDEM for timestep k.\\
$\Delta t_k$ & Interval length from the previous sample to the current sample.\\
$C_k$ & Set of collisions present in the current surfSurf sample.\\
$N_k = |C_k|$ & Number of sampled collisions at timestep k.\\
$N_k^{new}$ & Number of sampled collisions whose start time lies inside the current interval.\\
$s_i, e_i$ & Start time and end time of collision i reported by EDEM.\\
$\tau_i^{overlap}$ & Overlap duration of collision i with the interval $[t_{k-1}, t_k]$.\\
$d_i$ & Full contact duration of collision i.\\
$I[\cdot]$ & Indicator function equal to 1 when the stated condition is true and 0 otherwise.\\
$m_{1,i}, m_{2,i}$ & Masses of the two particles in collision i.\\
$r_{1,i}, r_{2,i}$ & Radii of the two particles in collision i.\\
$\mathbf{F}_{n,i}, \mathbf{F}_{s,i}$ & Current total normal and total tangential force vectors from EDEM.\\
$F_{n,i}, F_{s,i}, F_i$ & Magnitudes of normal, tangential, and total collision force used in the code.\\
$F_i^{max}$ & Maximum total collision-force magnitude formed from stored peak-force vectors.\\
$E_{n,i}, E_{s,i}, E_i$ & Normal, shear, and total collision energies used in the code.\\
$m_{eff,i}$ & Reduced mass of collision i.\\
$J_{n,i}, J_{s,i}, J_i$ & Normal, shear, and total impulse attributed to the current interval.\\
$A_i$ & Collision area proxy used by the exported stress calculation.\\
$\sigma_{n,i}, \sigma_{s,i}, \sigma_i$ & Exported normal, shear, and total stress-like values per collision.\\
$P_{n,i}, P_{s,i}, P_i$ & Interval-average power values per collision.\\
$n_{h,k}$ & Number of current sampled collisions touching host particle h at timestep k.\\
$\boldsymbol{\omega}_h$ & Angular-velocity vector of host particle h.\\
$d_{h,k}^{rot}, D_{h,k}^{rot}$ & Current and cumulative rotational travel distance for host particle h.\\
$\xi_{h,k}$ & Binary indicator showing whether host particle h is contacting at timestep k.\\
$T_{h,k}^{contact}$ & Cumulative time for which host particle h is marked as contacting.\\
$H_k$ & Host particles selected by the configured host material.\\
\bottomrule
\end{longtable}

\section*{How to read the symbols}
\paragraph{Time and interval symbols}
$k$ is a sample index and $t_k$ is the physical simulation time at that sample.
$\Delta t_k$ is the reporting interval width, so it is the denominator for rate-type quantities such as collision frequency and power.
Each collision has a start time $s_i$ and end time $e_i$, but only the overlap part $\tau_i^{overlap}$ is assigned to the current interval $[t_{k-1}, t_k]$.
By contrast, $d_i$ is the full lifetime of collision $i$ over its entire contact history, so $\tau_i^{overlap}$ and $d_i$ are not interchangeable.

\paragraph{Mechanical symbols}
$m_{1,i}, m_{2,i}$ and $r_{1,i}, r_{2,i}$ are the masses and radii of the two particles forming collision $i$.
Bold symbols $\mathbf{F}_{n,i}$ and $\mathbf{F}_{s,i}$ are vectors from EDEM, whereas $F_{n,i}$, $F_{s,i}$, and $F_i$ are magnitudes used by the export code.
$E$ terms denote energies, $m_{eff,i}$ is the reduced mass that normalizes those energies into specific intensity, and $J$ terms are impulses obtained by integrating force over time.

\paragraph{Stress and power symbols}
$A_i$ is a geometric area proxy built from particle radii, not a true Hertzian contact patch.
The exported stress-like values $\sigma_{n,i}$, $\sigma_{s,i}$, and $\sigma_i$ come from impulse divided by overlap time and by this area proxy.
They should be read as code-defined proxies, not as direct continuum stress or a true Hertz contact stress.
$P_{n,i}$, $P_{s,i}$, and $P_i$ are interval-average powers computed by dividing collision energy by $\Delta t_k$.

\paragraph{Host-particle rotation symbols}
For a host particle $h$, $\boldsymbol{\omega}_h$ is angular velocity, $n_{h,k}$ is the number of active sampled contacts touching that host particle, and $\xi_{h,k} = I[n_{h,k} > 0]$ converts that contact state into a binary gate.
That gate decides whether the current rotational travel distance $d_{h,k}^{rot}$ contributes to the cumulative rotational distance $D_{h,k}^{rot}$ and cumulative contact time $T_{h,k}^{contact}$.

\section*{Core equations}
\paragraph{C1. Interval length}
\[
\Delta t_k = \max(0, t_k - t_{k-1})
\]
For the first processed timestep, the implementation uses 0 as the previous time if no earlier sample is available.

\paragraph{C2. Overlap duration inside the current interval}
\[
\tau_i^{overlap} = \max\!\left(0, \min(e_i, t_k) - \max(s_i, t_{k-1})\right)
\]
This is the contact time attributed to the current interval for collision i.
It is not the same as the full contact duration $d_i$, which is defined separately in C3.

\paragraph{C3. Full contact duration of one collision}
\[
d_i = \max(0, e_i - s_i)
\]
Used for the summary field \code{Ave_contact_duration_per_collision}.

\paragraph{C4. Collision frequency based on new collisions started in the interval}
\[
\begin{aligned}
N_k^{new} &=
\begin{cases}
\sum_{i \in C_k} I[t_{k-1} \le s_i \le t_k], & t_{k-1} \le 0, \\
\sum_{i \in C_k} I[t_{k-1} < s_i \le t_k], & t_{k-1} > 0,
\end{cases} \\
f_k^{coll} &= \frac{N_k^{new}}{\Delta t_k}
\end{aligned}
\]
$C_k$ is the full set of collisions present in the current sample, not only the newly started ones. Therefore $N_k^{new}$ is only the subset count inside $C_k$ whose start times fall in the current interval. It is computed exactly as in \code{count_started_collisions} in the code: the first interval includes the lower bound so collisions starting exactly at time $0$ are counted once, while later intervals use an open lower bound to avoid double counting collisions that start exactly on a timestep boundary.

\paragraph{C5. Current force magnitudes}
\[
\begin{aligned}
F_{n,i} &= \|\mathbf{F}_{n,i}\|, \\
F_{s,i} &= \|\mathbf{F}_{s,i}\|, \\
F_i &= \|\mathbf{F}_{n,i} + \mathbf{F}_{s,i}\|
\end{aligned}
\]
The total collision force magnitude is the norm of the vector sum, not the sum of norms.

\paragraph{C6. Maximum-force magnitude used for impact statistics}
\[
F_i^{max} = \|\mathbf{F}_{n,i}^{max} + \mathbf{F}_{s,i}^{max}\|
\]
Used for \code{Ave_Fcol_Max_per_collision}, \code{Max_Fcol_Max_per_collision}, and \code{high_impact_event_fraction}.

\paragraph{C7. Reduced mass}
\[
m_{eff,i} = \frac{m_{1,i} m_{2,i}}{m_{1,i} + m_{2,i}}
\]
If the denominator is zero, the code treats the reduced mass as invalid and writes zero for specific-intensity terms.

\paragraph{C8. Collision specific-intensity values}
\[
\begin{aligned}
SI_{n,i} &= \frac{E_{n,i}}{m_{eff,i}}, \\
SI_{s,i} &= \frac{E_{s,i}}{m_{eff,i}}, \\
SI_i &= \frac{E_i}{m_{eff,i}}
\end{aligned}
\]
Here $E_i = E_{n,i} + E_{s,i}$. The current code sets dissipation energy to zero.

\paragraph{C9. Collision area proxy used in stress calculations}
\[
A_i = \pi \left(\frac{r_{1,i} + r_{2,i}}{2}\right)^2
\]
This is a geometric proxy area from the current implementation, not a Hertz contact patch.

\paragraph{C10. Trapezoidal impulse for matched sampled collisions}
\[
\begin{aligned}
J_{n,i} &= \frac{1}{2}\left(F_{n,i}^{prev} + F_{n,i}^{curr}\right)\tau_i^{overlap}, \\
J_{s,i} &= \frac{1}{2}\left(F_{s,i}^{prev} + F_{s,i}^{curr}\right)\tau_i^{overlap}, \\
J_i &= \frac{1}{2}\left(F_i^{prev} + F_i^{curr}\right)\tau_i^{overlap}
\end{aligned}
\]
The code matches current and previous collisions by unordered particle pair and, when needed, nearest contact position.
Only collisions that are present in the current sampled set and successfully matched to previous state contribute to this trapezoidal term; collision-level exports do not add a separate previous-only tail term.

\paragraph{C11. Collision stress values from impulse, overlap time, and area}
\[
\begin{aligned}
\sigma_{n,i} &= \frac{J_{n,i}}{\tau_i^{overlap} A_i}, \\
\sigma_{s,i} &= \frac{J_{s,i}}{\tau_i^{overlap} A_i}, \\
\sigma_i &= \frac{J_i}{\tau_i^{overlap} A_i}
\end{aligned}
\]
If overlap time or area is non-positive, the code writes zero.
These exported $\sigma$ values are therefore stress-like proxies from the implementation, not direct physical contact stresses.

\paragraph{C12. Collision power values}
\[
\begin{aligned}
P_{n,i} &= \frac{E_{n,i}}{\Delta t_k}, \\
P_{s,i} &= \frac{E_{s,i}}{\Delta t_k}, \\
P_i &= \frac{E_i}{\Delta t_k}
\end{aligned}
\]
Power uses the current interval length in the denominator.

\paragraph{C13. Aggregate specific energy by collision}
\[
\begin{aligned}
SE_{n,k}^{coll} &= \frac{\sum_{i \in C_k} E_{n,i}}{\sum_{i \in C_k} m_{eff,i}}, \\
SE_{s,k}^{coll} &= \frac{\sum_{i \in C_k} E_{s,i}}{\sum_{i \in C_k} m_{eff,i}}, \\
SE_k^{coll} &= \frac{\sum_{i \in C_k} E_i}{\sum_{i \in C_k} m_{eff,i}}
\end{aligned}
\]
The implementation aggregates energy first and then divides by the sum of valid reduced masses.
So this quantity is not the arithmetic mean of per-collision $SI_{n,i}$, $SI_{s,i}$, or $SI_i$ values.

\paragraph{C14. Particle split used for particle SI values still exported in the collision summary CSV}
\[
\begin{aligned}
E_{n,p} &= \frac{1}{2}\sum_{i \sim p} E_{n,i}, \\
E_{s,p} &= \frac{1}{2}\sum_{i \sim p} E_{s,i}, \\
SI_{n,p} &= \frac{E_{n,p}}{m_p}, \\
SI_{s,p} &= \frac{E_{s,p}}{m_p}
\end{aligned}
\]
These particle-based SI values appear in \code{collision_over_time.csv} because the shared timestep core computes both collision and particle summaries.

\paragraph{C15. Host-particle rotational metrics}
\[
\begin{aligned}
d_{h,k}^{rot} &= r_h \|\boldsymbol{\omega}_h\| \Delta t_k, \\
\xi_{h,k} &= I[n_{h,k} > 0], \\
D_{h,k}^{rot} &= D_{h,k-1}^{rot} + \xi_{h,k} d_{h,k}^{rot}, \\
T_{h,k}^{contact} &= T_{h,k-1}^{contact} + \xi_{h,k}\Delta t_k
\end{aligned}
\]
The contact indicator uses the current sampled active collision count for each host particle.
It is a gate based only on whether the host particle has any current sampled collision in $C_k$; it does not integrate overlap history by itself.

\section*{Summary CSV fields}
\small
\begin{longtable}{@{}L{0.33\textwidth}L{0.22\textwidth}L{0.37\textwidth}@{}}
\toprule
\textbf{Field} & \textbf{Calculation / reference} & \textbf{Meaning}\\
\midrule
\endhead
\code{time} & Current sample time $t_k$. & EDEM time value for the processed timestep.\\
\code{number_of_collisions} & $N_k = |C_k|$. & Number of sampled collisions in the current surfSurf snapshot.\\
\code{Ave_contact_duration_per_collision} & Mean of $d_i$ over current sampled collisions. Eq.~C3. & Average full contact duration.\\
\code{total_collision_time} & Sum of $\tau_i^{overlap}$ over sampled collisions. Eq.~C2. & Total overlap time attributed to the current interval.\\
\code{average_collision_time} & Mean of $\tau_i^{overlap}$ over sampled collisions. Eq.~C2. & Average overlap time per sampled collision in the current interval.\\
\code{collision_frequency} & $N_k^{new} / \Delta t_k$. Eq.~C4. & Frequency of the newly started subset of $C_k$, where $C_k$ itself is the full current sampled collision set.\\
\code{Ave_TotalColForce_inSys_byCollision} & Mean of $F_i$ over sampled collisions. Eq.~C5. & Average current total collision force magnitude.\\
\code{Ave_Fcol_Max_per_collision} & Mean of $F_i^{max}$. Eq.~C6. & Average maximum force magnitude.\\
\code{Max_Fcol_Max_per_collision} & Max of $F_i^{max}$. Eq.~C6. & Largest maximum force magnitude in the sampled collision set.\\
\code{high_impact_event_fraction} & Count($F_i^{max} > \mathrm{threshold}$) / $N_k$. Eq.~C6. & Fraction of sampled collisions above the configured impact threshold.\\
\code{Ave_SI_normal_per_collision} & Mean of $SI_{n,i}$ over valid reduced masses. Eq.~C8. & Average collision normal specific intensity.\\
\code{Ave_SI_shear_per_collision} & Mean of $SI_{s,i}$ over valid reduced masses. Eq.~C8. & Average collision shear specific intensity.\\
\code{Ave_SI_total_per_collision} & Mean of $SI_i$ over valid reduced masses. Eq.~C8. & Average collision total specific intensity.\\
\code{Total_normal_specific_energy_by_collision} & Aggregate collision normal specific energy. Eq. C13. & Energy summed over collisions, then divided by summed reduced mass.\\
\code{Total_shear_specific_energy_by_collision} & Aggregate collision shear specific energy. Eq. C13. & Energy summed over collisions, then divided by summed reduced mass.\\
\code{Total_specific_energy_by_collision} & Aggregate collision total specific energy. Eq. C13. & Energy summed over collisions, then divided by summed reduced mass.\\
\code{Ave_SI_normal_per_particle} & Mean over particles of $SI_{n,p}$ from the shared core. Eq.~C14. & Particle-based normal SI still written into the collision summary CSV.\\
\code{Ave_SI_shear_per_particle} & Mean over particles of $SI_{s,p}$ from the shared core. Eq.~C14. & Particle-based shear SI still written into the collision summary CSV.\\
\code{Ave_SI_total_per_particle} & Mean over particles of $SI_p$ from the shared core. Eq.~C14. & Particle-based total SI still written into the collision summary CSV.\\
\code{Ave_normal_stress_value_per_collision} & Mean of $\sigma_{n,i}$ over sampled collisions. Eq.~C11. & Average normal stress proxy per collision.\\
\code{Ave_shear_stress_value_per_collision} & Mean of $\sigma_{s,i}$ over sampled collisions. Eq.~C11. & Average shear stress proxy per collision.\\
\code{Ave_stress_value_per_collision} & Mean of $\sigma_i$ over sampled collisions. Eq.~C11. & Average total stress proxy per collision.\\
\code{Ave_normal_power_per_collision} & Mean of $P_{n,i}$ over sampled collisions. Eq.~C12. & Average normal power per collision.\\
\code{Ave_shear_power_per_collision} & Mean of $P_{s,i}$ over sampled collisions. Eq.~C12. & Average shear power per collision.\\
\code{Ave_power_per_collision} & Mean of $P_i$ over sampled collisions. Eq.~C12. & Average total power per collision.\\
\code{Total_normal_power_by_collision} & Sum of $P_{n,i}$ over sampled collisions. Eq.~C12. & Total normal power across the current sampled collision set.\\
\code{Total_shear_power_by_collision} & Sum of $P_{s,i}$ over sampled collisions. Eq.~C12. & Total shear power across the current sampled collision set.\\
\code{Total_power_by_collision} & Sum of $P_i$ over sampled collisions. Eq.~C12. & Total power across the current sampled collision set.\\
\code{Ave_rotational_distance_per_host_particle} & Mean of $d_{h,k}^{rot}$ over all host particles sampled at $k$. Eq.~C15. & Average host rotational distance in the current interval.\\
\code{Ave_rotational_distance_total_per_host_particle} & Mean of $D_{h,k}^{rot}$ over host particles with $\xi_{h,k} = 1$. Eq.~C15. & Average cumulative host rotational distance among currently contacting host particles.\\
\code{Ave_cumulative_contact_time_per_host_particle} & Mean of $T_{h,k}^{contact}$ over host particles with $\xi_{h,k} = 1$. Eq.~C15. & Average cumulative host contact time among currently contacting host particles.\\
\bottomrule
\end{longtable}
\normalsize

\section*{Extracted CSV fields}
\small
\begin{longtable}{@{}L{0.33\textwidth}L{0.22\textwidth}L{0.37\textwidth}@{}}
\toprule
\textbf{Field} & \textbf{Calculation / reference} & \textbf{Meaning}\\
\midrule
\endhead
\code{time} & Scalar $t_k$ for the extract timestep. & Written once at the top of the extracted block by the CSV expansion helper.\\
\code{collision_id} & Row index $0, 1, \ldots, N_k - 1$. & Local collision row number in the current extract snapshot, not a persistent collision ID.\\
\code{first_particle_id} & Raw first collision particle ID from EDEM. & First particle in the sampled collision pair.\\
\code{second_particle_id} & Raw second collision particle ID from EDEM. & Second particle in the sampled collision pair.\\
\code{specific_energy_normal_per_collision} & $SI_{n,i}$ when reduced mass is valid, else $0$. Eq.~C8. & Collision normal specific intensity.\\
\code{specific_energy_shear_per_collision} & $SI_{s,i}$ when reduced mass is valid, else $0$. Eq.~C8. & Collision shear specific intensity.\\
\code{specific_energy_total_per_collision} & $SI_i$ when reduced mass is valid, else $0$. Eq.~C8. & Collision total specific intensity.\\
\code{collision_area} & $A_i$ from Eq.~C9. & Collision area proxy used by the stress calculation.\\
\code{normal_force_magnitude_per_collision} & $F_{n,i}$ from Eq.~C5. & Current normal-force magnitude.\\
\code{shear_force_magnitude_per_collision} & $F_{s,i}$ from Eq.~C5. & Current tangential-force magnitude.\\
\code{total_force_magnitude_per_collision} & $F_i$ from Eq.~C5. & Current total collision-force magnitude.\\
\code{normal_stress_value_per_collision} & $\sigma_{n,i}$ from Eq.~C11. & Normal stress proxy.\\
\code{shear_stress_value_per_collision} & $\sigma_{s,i}$ from Eq.~C11. & Shear stress proxy.\\
\code{stress_value_per_collision} & $\sigma_i$ from Eq.~C11. & Total stress proxy.\\
\code{normal_power_per_collision} & $P_{n,i}$ from Eq.~C12. & Normal power.\\
\code{shear_power_per_collision} & $P_{s,i}$ from Eq.~C12. & Shear power.\\
\code{power_per_collision} & $P_i$ from Eq.~C12. & Total power.\\
\bottomrule
\end{longtable}
\normalsize

\section*{Implementation notes}
\noindent - All displayed formulas reflect the current implementation in \code{stress_collision_based.py}.\par
\noindent - Collision continuity for the trapezoidal rule is matched by unordered particle pair and nearest contact position when multiple contacts share the same pair.\par
\noindent - The collision summary and collision extracted file do not add a separate previous-only tail term. Tail handling is only added at particle level inside the shared timestep core.\par
\noindent - If a denominator is zero, if no valid term exists, or if a feature flag is disabled in \code{Stress_settings.txt}, the corresponding exported values remain 0.\par
\noindent - The CSV expansion helper writes scalar fields such as time only once at the top of an extracted block, leaving the remaining rows blank in that scalar column.\par

\end{document}
